\subsection{sepal}

\begin{frame}
	System for Earth Observation Data Access, Processing and Analysis for Land
	Monitoring - SEPAL.
\end{frame}

\begin{frame}
	\frametitle{SEPAL}
	\begin{columns}
		\begin{column}{0.5\textwidth}
			\begin{itemize}
				\item Cloud computing-based platform for autonomous land
				      monitoring using remote sensing data.
				\item This platform allows users to access powerful
				      cloud-computing resources to query, access and process
				      satelllite data quickly and efficiently for conducting
				      advanced analysis.
				\item It is part of OpenForis.
			\end{itemize}
		\end{column}
		\begin{column}{0.5\textwidth}
			\begin{figure}
				\centering
				\includegraphics[scale=0.4]{logos/sepal.png}
				\includegraphics[scale=0.8]{logos/openforis.png}
			\end{figure}
		\end{column}
	\end{columns}
\end{frame}

\begin{frame}
	\frametitle{Setup}
	Create accounts at:
	\begin{itemize}
		\item SEPAL.
		\item Google Earth Engine.
		\item Collect Earth Engine.
	\end{itemize}
\end{frame}

\begin{frame}
	\frametitle{Ask for credit}
	\begin{itemize}
		\item SEPAL runs on top of the cloud computing platform provided by
		      Amazon.
		\item FAO pays for your use of SEPAL resources as long as you use them
		      to process Earth Observation data.
		\item Once you open your account, you have no credit and you need to
		      ask for it. Check the instructions %\href{https://docs.sepal.io/en/latest/setup/register.html#request}{here}.
	\end{itemize}
\end{frame}

\begin{frame}
	\frametitle{Get data}
	\begin{itemize}
		\item Connect SEPAL to Google Earth Engine.
		      %\item Connect SEPAL to NICFI-PlanetLab (Norway's International Climate 
		      %    and Forest Initiative).
	\end{itemize}
\end{frame}

\begin{frame}
	\frametitle{sepal.io}
	The platform is organized into:
	\begin{itemize}
		\item Recipes.
		\item Modules.
		\item Workflows.
	\end{itemize}
\end{frame}

\begin{frame}
	\frametitle{Recipes}
	\begin{itemize}
		\item Quickly and efficiently query and process satellite data.
		\item A recipe is a record of steps and parameters used to make a
		      data set (e.g. classification).
		\item Access GEE imagery catalog and run planetary-scale analysis
		      without a line of code.
		\item Recipes include creating mosaics (optical, radar, planet),
		      supervised classifications (of images or time series).
	\end{itemize}
\end{frame}

\begin{frame}
	\frametitle{Modules}
	\begin{itemize}
		\item GIS tools that complement SEPAL's recipes.
		\item Based on advanced GIS libraries.
	\end{itemize}
\end{frame}

\begin{frame}
	\frametitle{Workflows}
	\begin{itemize}
		\item Combinations of Recipes, modules, and tools to perform complex
		      data analysis.
	\end{itemize}
\end{frame}

\begin{frame}
	\frametitle{CLI utilities}
	\begin{itemize}
		\item Command Line Interface utilities.
		\item GDAL, Google Drive, GEE, GuidosToolbox Workbench, Open Foris
		      Geospatial Toolbox, Orfeo Toolbox, Python, R code.
		\item IDEs.
	\end{itemize}
\end{frame}

\begin{frame}
	\frametitle{IDEs}
	\begin{itemize}
		\item Integrated Development Environments.
		\item JupyterLab.
		\item Jupyter Notebook.
		\item RStudio.
	\end{itemize}
\end{frame}

\begin{frame}
	\frametitle{Start a machine}
\end{frame}

\begin{frame}
	\frametitle{Get data}
\end{frame}

\begin{frame}
	\frametitle{Start an R-spatial container}
\end{frame}

\begin{frame}
	\frametitle{Get root password}
\end{frame}
